% Exemplo: quadro com cabeçalhos centralizados (multicolumn de uma coluna) e letra reduzida
% Compilar: latexmk -pdf quadro-tecnologias.tex   Testado com: abntex2, TeX Live 2024 ou mais recente
\documentclass[12pt, a4paper, brazilian]{abntex2}
\usepackage[T1]{fontenc}
\usepackage{tabularx}

%%BEGIN:quadropreambulo
% Float "Quadro" (a abntex2 não o define): numeração própria e lista de quadros (memoir)
\newfloat{quadro}{loq}{Quadro}
\newlistof{listofquadros}{loq}{Lista de quadros}
\newlistentry{quadro}{loq}{0}         % linha da lista, com pontilhado e página
\providecommand{\quadroautorefname}{Quadro}  % nome usado pelo \autoref
%%END:quadropreambulo

\begin{document}

As tecnologias estão no \autoref{qua:tecnologias}.

%%BEGIN:quadrotecnologias
\begin{quadro}[htbp]
  \caption{Tecnologias empregadas e justificativas (exemplo)}\label{qua:tecnologias}
  \small
  \renewcommand{\arraystretch}{1.25}
  \begin{tabularx}{\textwidth}{|l|l|X|}
    \hline
    \multicolumn{1}{|c|}{\textbf{Camada}} & \multicolumn{1}{c|}{\textbf{Tecnologia}} & \multicolumn{1}{c|}{\textbf{Justificativa}} \\
    \hline
    Interface & React 18 & Componentização e grande comunidade; facilita a manutenção da interface. \\
    \hline
    Serviços & Python 3 e FastAPI & Produtividade, documentação automática da API e desempenho adequado ao porte do sistema. \\
    \hline
    Dados & PostgreSQL & SGBD relacional de código aberto, com transações e suporte a JSON. \\
    \hline
  \end{tabularx}
  \fonte{Elaborado pelo próprio autor.}
\end{quadro}
%%END:quadrotecnologias

\end{document}
